% texte.tex — f164r, vv. 123–164 (conversion initiale des travaux antérieurs)
\vers[n=123, indent=4em, findent=2.4em]{}{}{σέβομαι Λατώ τ' ἄνασσαν}{8° — J'honore Léto la souveraine}
\vers[n=124, indent=4em, findent=2.4em]{}{}{κίθαρίν τε ματέρ᾽ ὕμνων}{et la cithare mère des hymnes,}
\vers[n=125, indent=4em, findent=2.4em]{}{}{ἄρσενι βοᾷ δοκιμῳ.}{à la voix mâle et éprouvée.}
\vers[n=126, indent=4em, findent=2.4em]{}{}{τᾷ φάος ἔσσυτο δαιμονίοις}{9° — À ces accents, une lumière jaillit}
\vers[n=126, indent=4em, findent=2.4em]{}{}{ὄμμασιν ἁμετέρας τε δι᾽ αἰφνι-}{dans les yeux des dieux}
\vers[suite]{ΑΓ.}{Agathon}{\alinea δίου ὀπός. ὧν χάριν}{et traverse notre soudain regard.}
\vers[n=128, indent=4em, findent=2.4em]{}{}{ἄνακτ᾽ ἄγαλλε Φοῖβον.}{Aussi bien, glorifie le seigneur Phébus.}
\vers[n=129, indent=4em, findent=2.4em]{}{}{χαῖρ᾽ ὄλβιε παῖ Λατοῦς.}{10° — Salut, heureux fils de Léto.}
\vers[n=130]{ΜΝ.}{Mnésiloque}{ὡς ἡδὺ τὸ μέλος ὦ πότνιαι Γενετυλλίδες}{Quel chant suave, ô Génétyllides vénérées}
\vers[n=131]{}{}{καὶ θηλυδριῶδες καὶ κατεγλωττισμένον}{et qui sent la féminie et les coups de langue}
\vers[n=132]{}{}{καὶ μανδαλωτόν, ὥστ᾽ ἐμοῦ γ᾽ ἀκροωμένου}{et les baisers lascifs. En l'écoutant}
\vers[n=133]{}{}{ὑπὸ τὴν ἕδραν αὐτὴν ὑπῆλθε γάργαλος.}{j'ai ressenti un prurit jusque sous le fondement.}
\vers[n=134]{}{}{καί σ᾽ ὦ νεανίσκ᾽ εἴ τις εἶ, κατ᾽ Αἰσχύλον}{Jouvenceau, si tu en es un, à la façon d'Eschyle}
\vers[n=135]{}{}{ἐκ τῆς Λυκουργείας ἐρέσθαι βούλομαι.}{dans sa \italique{Lycurgie}, je veux t'interroger.}
\vers[n=136]{}{}{ποδαπὸς ὁ γύννις; τίς πάτρα; τίς ἡ στολή;}{D'où es-tu, ô femmelette ? Quelle est ta patrie ? Qu'est cet accoutrement,}
\vers[n=137]{}{}{τίς ἡ τάραξις τοῦ βίου; τί βάρβιτος}{cette incohérence dans tes façons ? Un barbitos,}
\vers[n=138]{}{}{λαλεῖ κροκωτῷ; τί δὲ λύρα κεκρυφάλῳ;}{qu'a-t-il à dire à une crocote, une lyre à une résille ?}
\vers[n=139]{}{}{τί λήκυθος καὶ στρόφιον; ὡς οὐ ξύμφορον.}{Quoi ! une pomme d'Adam avec un strophion ? Quelle disconvenance !}
\vers[n=140]{}{}{τίς δαὶ κατρόπτου καὶ ξίφους κοινωνία;}{Le rapport d'un miroir à une épée ?}
\vers[n=141]{}{}{σύ τ᾽ αὐτὸς ὦ παῖ; πότερον ὡς ἀνὴρ τρέφει;}{Et toi-même, ô jeunesse, qui es-tu ? Un homme ?}
\vers[n=142]{}{}{καὶ ποῦ πέος; ποῦ χλαῖνα; ποῦ Λακωνικαί;}{Mais où est ton cas ? où ton manteau ? où tes laconiennes ?}
\vers[n=143]{}{}{ἀλλ᾽ ὡς γυνὴ δῆτ᾽ εἶτα ποῦ τὰ τιτθία;}{Alors tu es femme ? Mais où sont tes seins ?}
\vers[n=144]{}{}{τί φῄς; τί σιγᾷς; ἀλλὰ δῆτ᾽ ἐκ τοῦ μέλους}{Que dis-tu ? Pourquoi te taire ? Alors d'après ton chant,}
\vers[n=145]{}{}{ζητῶ σ᾽, ἐπειδή γ᾽ αὐτὸς οὐ βούλει φράσαι;}{je te déchiffre puisque tu ne veux pas t'expliquer toi-même.}
\vers[n=146]{ΑΓ.}{Agathon}{ὦ πρέσβυ πρέσβυ, τοῦ φθόνου μὲν τὸν ψόγον}{Vieillard, vieillard, ces reproches dictés par l'envie,}
\vers[n=147]{}{}{ἤκουσα, τὴν δ᾽ ἄλγησιν οὐ παρεσχόμην‎}{je les ai entendus, mais n'en ai point senti l'atteinte.}
\vers[n=148]{}{}{ἐγὼ δὲ τὴν ἐσθῆθ᾽ ἅμα τῇ γνώμῃ φορῶ.}{Je porte la mise congruente à mes pensées.}
\vers[n=149]{}{}{χρὴ γὰρ ποιητὴν ἄνδρα πρὸς τὰ δράματα}{Car c'est une loi que le poète, suivant les drames}
\vers[n=150]{}{}{ἃ δεῖ ποιεῖν πρὸς ταῦτα τοὺς τρόπους ἔχειν.}{qu'il lui faut faire, modèle sur eux ses mœurs.}
\vers[n=151]{}{}{αὐτίκα γυναικεῖ᾽ ἢν ποιῇ τις δράματα,}{Et par exemple, si l'on compose des drames féminins,}
\vers[n=152]{}{}{μετουσίαν δεῖ τῶν τρόπων τὸ σῶμ᾽ ἔχειν.}{il faut que le corps participe aux mœurs de la femme..}
\vers[n=153]{ΜΝ.}{Mnésiloque}{οὐκοῦν κελητίζεις, ὅταν Φαίδραν ποιῇς;}{Ainsi tu chevauches quand tu composes une Phèdre ?}
\vers[n=154]{ΑΓ.}{Agathon}{ἀνδρεῖα δ᾽ ἄν ποιῇ τις, ἐν τῷ σώματι}{Compose-t-on des drames mâles, dans le corps}
\vers[n=155]{}{}{ἔνεσθ᾽ ὑπάρχον τοῦθ᾽. ἃ δ᾽ οὐ κεκτήμεθα,}{réside ce qui est requis. Mais ce que nous ne possédons pas,}
\vers[n=156]{}{}{μίμησις ἤδη ταῦτα συνθηρεύεται.}{l'imitation est là qui l'attrape.}
\vers[n=157]{ΜΝ.}{Mnésiloque}{ὅταν σατύρους τοίνυν ποιῇς, καλεῖν ἐμέ,}{Eh bien, quand tu feras des drames satyriques, appelle-moi,}
\vers[n=158]{}{}{ἵνα συμποιῶ σοὔπισθεν ἐστυκὼς ἐγώ.}{et je coopérerai avec toi, de derrière, avec raideur.}
\vers[n=159]{ΑΓ.}{Agathon}{ἄλλως τ᾽ ἄμουσόν ἐστι ποιητὴν ἰδεῖν}{D'ailleurs, c'est chose choquante à voir qu'un poète}
\vers[n=160]{}{}{ἀγρεῖον ὄντα καὶ δασύν σκέψαι δ᾽ ὅτι}{grossier et hirsute. Avise-toi que}
\vers[n=161]{}{}{Ἴβυκος ἐκεῖνος κἀνακρέων ὁ Τήιος}{le fameux Ibycos, et Anacréon de Téos}
\vers[n=162]{}{}{κἀλκαῖος, οἳ περὶ ἁρμονίαν ἐκεχήνεσαν,}{et Alcée, accordaient leur voix,}
\vers[n=163]{}{}{ἐμιτροφόρουν τε καὶ διέκειντ' Ἰωνικῶς,}{leur mitre et leur vit à l'ionienne. Et Phrynichos,}
\vers[n=164]{}{}{καὶ Φρύνιχος, τοῦτον γὰρ οὖν ἀκήκοας,}{— celui-là, en effet, tu l'as entendu, —}
